\documentclass[10pt,letterpaper]{article}
\usepackage[top=0.85in,left=2.75in,footskip=0.75in]{geometry}

% amsmath and amssymb packages, useful for mathematical formulas and symbols
\usepackage{amsmath,amssymb}

% Use adjustwidth environment to exceed column width (see example table in text)
\usepackage{changepage}

% Use Unicode characters when possible
\usepackage[utf8x]{inputenc}

% textcomp package and marvosym package for additional characters
\usepackage{textcomp,marvosym}

% cite package, to clean up citations in the main text. Do not remove.
\usepackage{cite}

% Use nameref to cite supporting information files (see Supporting Information section for more info)
\usepackage{nameref,hyperref}

% line numbers
\usepackage[right]{lineno}

% ligatures disabled
\usepackage{microtype}
\DisableLigatures[f]{encoding = *, family = * }

% color can be used to apply background shading to table cells only
\usepackage[table]{xcolor}

% array package and thick rules for tables
\usepackage{array}

% enumerate package lets us use letters instead of numbers
\usepackage{enumerate}

% create "+" rule type for thick vertical lines
\newcolumntype{+}{!{\vrule width 2pt}}

% create \thickcline for thick horizontal lines of variable length
\newlength\savedwidth
\newcommand\thickcline[1]{%
  \noalign{\global\savedwidth\arrayrulewidth\global\arrayrulewidth 2pt}%
  \cline{#1}%
  \noalign{\vskip\arrayrulewidth}%
  \noalign{\global\arrayrulewidth\savedwidth}%
}

% \thickhline command for thick horizontal lines that span the table
\newcommand\thickhline{\noalign{\global\savedwidth\arrayrulewidth\global\arrayrulewidth 2pt}%
\hline
\noalign{\global\arrayrulewidth\savedwidth}}

\usepackage{color}

% Remove comment for double spacing
%\usepackage{setspace}
%\doublespacing

% Text layout
\raggedright
\setlength{\parindent}{0.5cm}
\textwidth 5.25in
\textheight 8.75in

% Bold the 'Figure #' in the caption and separate it from the title/caption with a period
% Captions will be left justified
\usepackage[aboveskip=1pt,labelfont=bf,labelsep=period,justification=raggedright,singlelinecheck=off]{caption}
\renewcommand{\figurename}{Fig}

% Use the PLoS provided BiBTeX style
\bibliographystyle{plos2015}

% Remove brackets from numbering in List of References
\makeatletter
\renewcommand{\@biblabel}[1]{\quad#1.}
\makeatother

% Leave date blank
\date{}

% Header and Footer with logo
\usepackage{lastpage,fancyhdr,graphicx}
\usepackage{epstopdf}
\pagestyle{myheadings}
\pagestyle{fancy}
\fancyhf{}
\setlength{\headheight}{27.023pt}
\lhead{\includegraphics[width=2.0in]{PLOS-submission.eps}}
\rhead{\Large{DRAFT}}
\rfoot{\thepage/\pageref{LastPage}}
\renewcommand{\footrule}{\hrule height 2pt \vspace{2mm}}
\fancyheadoffset[L]{2.25in}
\fancyfootoffset[L]{2.25in}
\lfoot{\sf PLOS}


% Define per-paper macros.
\newcommand{\withurl}[2]{{#1}\footnote{\texttt{#2}}}
\newcommand{\rulemajor}[1]{\section{#1}}
\begin{document}
\vspace*{0.2in}

\begin{flushleft}
{\Large
\textbf\newline{Ten Quick Tips for Running a Less Painful Meeting}
}
\newline
\\
{Greg~Wilson}\textsuperscript{1,*}
\\
\textbf{1} RStudio, Inc., Toronto, Ontario M4L 2T9
\\
\bigskip
* greg.wilson@rstudio.com
\end{flushleft}

\section*{Abstract}

Researchers complain that meetings are a waste of time, but that is partly
because most have never learned how to run one or participate in one
efficiently. The ten quick tips below will make meetings less frequent, shorter,
and more effective.

\section*{Author Summary}

Knowing how to run a meeting efficiently is a core skill for anyone who wants to
get things done, as is knowing how to participate in someone else's meeting.
This article gives ten quick tips for doing both.

\section*{Introduction}

Most people are really bad at meetings: they don't have an agenda going in, they
don't take minutes, they wander off into irrelevancies, they repeat what others
have said or recite banalities simply so that they'll have said something, and
they devote half of their attention to their email or to side conversations.
Knowing how to run a meeting efficiently is a core skill for anyone who wants to
get things done.  Knowing how to take part in someone else's meeting is just as
important, but gets far less attention. The most important rules for making
meetings efficient are not secret, but are rarely followed; this article gives
ten quick tips that will make meetings less frequent, shorter, and more
effective.

\rulemajor{Decide if there actually needs to be a meeting.}

If the only purpose is to share information, have everyone send a brief email
instead, or better yet, add notes to a shared document so that people can find
the information later. Remember, you can read faster than anyone can speak: if
someone has facts for the rest of the team to absorb, the most polite way to
communicate them is to type them in.  And if people don't read what you collate
and send out, you can always offer them the option of having a meeting instead.
In almost thirty years, I've never had anyone take me up on this{\ldots}

\rulemajor{Create an agenda.}

If nobody cares enough about the meeting to make a point-form list of what's
supposed to be discussed, or if that list isn't circulated at least a day in
advance, the meeting probably doesn't need to happen.  (``The agenda is
everything we've talked about recently on Slack'' doesn't count.)

\textbf{Include timings in the agenda.}  Timings help you keep the meeting
moving, as in, ``That's very interesting, but we have five other topics to get
through in the next fifteen minutes, so could you please make your point?''.
Your first estimates with any new group will be wildly optimistic, so revise
them upward for subsequent meetings.

\textbf{Prioritize.}  Every meeting is a micro-project, so work should be
prioritized in the same way that it is for other projects: things that will have
high impact but take little time should be done first, and things that will take
lots of time but have little impact should be skipped entirely.

\textbf{Plan to end early.}  If your meeting is scheduled for 10:00-11:00, you
should aim to end at 10:50 to give people time to get where they need to go
next.

\textbf{Don't include a quiz.}  There is a growing tendency in tech companies to
have people answer a question at the end of the agenda to prove that they've
read it.  Don't do this: one of the first people to look at the agenda will warn
everyone else, so all you'll do is make people feel that you're treating them
like children.

\rulemajor{Have clear rules for making decisions.}

If your meeting includes more than a couple of dozen people, or if you think you
need Robert's Rules of Order for any other reason, you should probably get some
training in how to run and participate in meetings.  For any smaller group, or
for groups in which people mostly agree with each other, use Martha's Rules
instead~\cite{Minahan1986}.

\begin{enumerate}

\item
  Before each meeting, anyone who wishes may sponsor a proposal.  Proposals must
  be circulated at least 24 hours before a meeting in order to be considered at
  that meeting, and must include:

  \begin{itemize}
  \item a one-line summary (the subject line of the issue);
  \item the full text of the proposal;
  \item any required background information;
  \item pros and cons; and
  \item possible alternatives.
  \end{itemize}

\item
  A quorum is established in a meeting if half or more of voting members are
  present.

\item
  Once a person has sponsored a proposal, they are responsible for it.  The
  group may not discuss or vote on the issue unless the sponsor or their
  delegate is present.  The sponsor is also responsible for presenting the item
  to the group.

\item
  After the sponsor presents the proposal, a \emph{sense vote} is cast for the
  proposal prior to any discussion:

  \begin{itemize}
  \item Who likes the proposal?
  \item Who can live with the proposal?
  \item Who is uncomfortable with the proposal?
  \end{itemize}

\item
  If all or most of the group likes or can live with the proposal, it is
  immediately moved to a formal vote with no further discussion.

\item
  If most of the group is uncomfortable with the proposal, it is postponed for
  further rework by the sponsor.

\item
  If some members are uncomfortable they can briefly state their objections.  A
  timer is then set for a brief discussion moderated by the facilitator.  After
  10 minutes or when no one has anything further to add (whichever comes first),
  the facilitator calls for a yes-or-no vote on the question: ``Should we
  implement this decision over the stated objections?''  If a majority votes
  "yes" the proposal is implemented.  Otherwise, the proposal is returned to the
  sponsor for further work.

\end{enumerate}

\rulemajor{Put someone in charge.}

``In charge'' means keeping the meeting moving, glaring at people who are
muttering to one another or checking email, and telling people who are talking
too much to get to the point.  It does \emph{not} mean doing all the talking; in
fact, whoever is in charge will usually talk less than anyone else, just as a
referee usually kicks the ball less often than the players.  One way to ensure
this is to give the moderator a stuffed animal or something else to hold up when
they're speaking on their own behalf rather than moderating.

If you are the moderator, \textbf{keep a backlog.}  If anything comes up that
looks like it will take the conversation into the weeds, write it down on a
sticky note and put it on the table (or the wall behind the moderator, or
wherever else is convenient), and then deal with it at the end of the meeting.
This works better than having the moderator take a note, since everyone can see
the backlog piling up and compare it against the time remaining.

Also, \textbf{call on specific people} to make sure everyone's ideas are heard.
Some people are more introverted than others, or need a little more time to
think, or have been trampled so often by fast talkers that they've given up
trying to speak. Keep a tally of who's spoken how often and try to draw out
those who are falling behind.

\rulemajor{Require politeness.}

No one gets to be rude, no one gets to ramble, and if someone goes off topic,
it's the chair's job to say, ``Let's discuss that elsewhere.''  If your project
has a code of conduct (which it should), the moderator should remind people of
it at the start of each meeting.

If attendees are co-located, enforce a \textbf{no technology} rule: everyone
must put their phones, tablets, and laptops into politeness mode (i.e., close
them).  If this is too stressful, let participants hang on to their electronic
pacifiers, but turn off the network so that they really \emph{are} using them
just to take notes or check the agenda.  Only allow exceptions if people need
their devices for accessibility reasons.

\textbf{No interruptions.}  Participants should raise a finger or make one of
the demure gestures used at high-priced auctions if they want to speak. It is
the moderator's job to keep track of the speaking queue, and to call on people
in order.

\rulemajor{Record minutes.}

Someone other than the chair should take point-form notes about the most
important pieces of information that were shared, about every decision that was
made, and about every task that was assigned to someone.  As soon as the meeting
is over, the minutes should be circulated (e.g., emailed to everyone or posted
to a wiki):

\begin{itemize}

\item
  \textbf{People who weren't at the meeting can keep track of what's going on.}
  You and your colleagues all have to juggle multiple tasks.  A wiki page, email
  message, or blog entry is more efficient and more accurate than asking a
  colleague what you missed.

\item
  \textbf{Everyone can check what was actually said or promised.} More than
  once, I've looked over the minutes of a meeting I was in and thought, ``I
  didn't promise to have it ready then!''  Accidentally or not, people will
  often remember things differently; writing it down gives team members a chance
  to correct mistaken or malicious interpretations.

\item
  \textbf{People can be held accountable at subsequent meetings.} There's no
  point assigning action items if you don't follow up on them.  If you're using
  a ticketing system such as GitHub Issues, the best thing to do is to create a
  ticket for each new question or task right after the meeting, and update those
  that are being carried forward.  The set of open issues can then serve as a
  starting point for the next meeting's agenda.

\end{itemize}
  
\rulemajor{Manage ``that guy''.}

Some people are so used to the sound of their own voice that they will insist on
talking half the time no matter how many other people are in the room. One way
to combat this is to give everyone three sticky notes at the start of the
meeting. Every time they speak, they have to take down one sticky note. When
they're out of notes, they aren't allowed to speak until everyone has used at
least one, at which point everyone gets all of their sticky notes back. This
ensures that nobody talks more than three times as often as the quietest person
in the meeting, and completely changes the dynamics of most groups: people who
have given up trying to be heard because they always get trampled suddenly have
space to contribute, and the overly-frequent speakers quickly realize just how
unfair they have been.

Another useful technique is ``interruption bingo''. Draw a grid and label the
rows and columns with the participants' names. Each time someone interrupts
someone else, add a tally mark to the appropriate cell. Halfway through the
meeting, take a moment to look at the results.  In most cases, you will see that
one or two people are doing all of the interrupting, often without being aware
of it. After that, saying, ``All right, I'm adding another tally to the bingo
card,'' is often enough to get them to throttle back.

\rulemajor{Be an active participant.}

Everyone offers leadership training, but nobody offers followership training.
The result is that many meetings are full of conductors who have no idea how to
actually play in an orchestra.  The rules for doing better are natural
complements of the ones given above.

\textbf{Decline invitations.} If you have read the agenda and don't think you
need to be in a meeting, say so, but only if you are also willing to agree in
advance to abide by the decisions made in the meeting.

This rule requires that you \textbf{read the agenda and supporting material
\emph{before} the meeting.}  If you still don't know after doing this whether
you need to attend or not, ask the moderator to clarify the agenda; this will
make the meeting more efficient for everyone.

\textbf{Take your own notes.}  While other people are talking, take notes of
questions you want to ask or points you want to make.  You'll be surprised how
smart it makes you look when it's your turn to speak.

\textbf{Use other participants' names.}  This is particularly important in
online meetings, since you can't use body language to indicate who you're
addressing.

\textbf{Pause before speaking.} If the moderator is managing the queue of
speakers so that you don't lose your chance to speak if you hesitate for a
couple of heartbeats, make it a habit to pause for a moment before speaking.  It
may not change what you say, but it's almost certain to make you sound more
eloquent.

\textbf{Put down your hand.} If someone else has already made the point you were
going to make, take yourself out of the speaking queue, or say, ``I agree
with{\ldots}'' and leave it at that.

\rulemajor{Modify these rules for online meetings.}

Depending on your workplace, many or most of your  meetings will take place via
video conferencing. Most of the previous rules apply directly to such meetings,
but a few modifications are needed.

First, meetings should be \textbf{all in person or all online}. If some people
are co-located and some are not, the remote participants are at a significant
disadvantage. The easiest way to level the playing field is to mandate that if
even one person is remote, everyone has to take part via video conferencing.

Second, \textbf{do \emph{not} record the meeting}: doing so can inhibit
participation (particularly from anyone who has ever been the victim of online
harassment), and people who say they want a recording in order to find out what
was said should be reading the minutes instead.

Third, since you've already broken the ``no technology'' rule, \textbf{take
minutes in a shared document} so that everyone can see them in real time.

Fourth, remind everyone that audio matters more than video, so they should use a
headset microphone rather than their phone or the microphone in their laptop.
Mute everyone at the start of the meeting, and re-mute people each time they
finish speaking.

Finally, Chelsea Troy's discussion of why online meetings are often frustrating
and unproductive \cite{Troy2018} makes an important point: in most online
meetings, the first person to speak during a pause gets the floor.  The result?
``If you have something you want to say, you have to stop listening to the
person currently speaking and instead focus on when they're gonna pause or
finish so you can leap into that nanosecond of silence and be the first to utter
something.  The format{\ldots}encourages participants who want to contribute to
say more and listen less.''  The solution is to run a text chat beside the video
conference where people can signal that they want to speak, The moderator then
selects people from the waiting list.

\rulemajor{Seek truth, not victory.}

The last rule is the hardest for many people, but also the most important.
Meetings can all too easily degenerate into struggles for social dominance, and
people can all too easily find themselves arguing to win rather than to find
solutions.  The moderator should keep this under control, but they can't do it
without participants' support.

In particular, \textbf{don't raise a point you don't actually agree with} just
to air it or because you haven't spoken in a while.  The devil doesn't actually
need more advocates, and even if he did, you don't have to be one.

\section*{Conclusion}

Different cultures have different conventions, and every group will evolve its
own practices for meetings as time goes on~\cite{Bollier2014}. If you would like
to accelerate that process, \cite{Brookfield2016} catalogs fifty different ways
to get groups talking productively, most of which were new to the author, and
only half a dozen of which have ever been supported directly online.

\bibliography{10-meeting}

\end{document}
